\documentclass{nasa-handbook} % [letter] for US Letter

\nhsetup{
  doc         = {[XX2A]-03-BLOCK I-(1)},
  title       = {[PROGRAM] OPERATIONS HANDBOOK},
  part        = {SYSTEMS DATA},
  section     = 2,
  subsection  = 2.1,
  subsystem   = {[SUBSYSTEM NAME]},
  tab         = {[TAB]},
  basic date  = {[1 January 1970]},
  change date = {},
  mission     = {},
  mission label = {Mission}, % System or Project for IT subjects
}

\begin{document}

\nhopener

\begin{nhpara}{2.1.1}{INTRODUCTION.}
[Three lines: what the subsystem is, the one idea the reader must leave with, and a pointer to the hero figure (figure 2.1-1).]
\end{nhpara}

% FIGURE 2.1-1: the hero, 66-72 mm tall. Units are mm, origin top left, y down,
% 180 wide. Arguments: [top y], height, number, caption, figure code.
\begin{nhfigure}{66}{2.1-1}{[Title Case Caption]}{[TAB]-501A}
  \path[panel] (1.5,2.5) rectangle (81.5,62.5);
  \node[lbl-s] at (3.4,5.6) {[PANEL NAME]};
  \path[box] (8,20) rectangle (34,31);
  \node[lbl,c] at (21,24.8) {[COMPONENT]};
  \node[lbl-s,c] at (21,27.8) {[ONE DETAIL]};
  \path[wa] (34,25.5) -- (55.6,25.5);
  \node[lbl-s] at (36,24.4) {[SIGNAL]};
  \path[box] (56,20) rectangle (78,31);
  \node[lbl,c] at (67,26.2) {[NEXT]};
\end{nhfigure}

\vspace{2.2mm}
\begin{nhpara}{2.1.2}{FUNCTIONAL DESCRIPTION.}
[Four lines: how it works, in the order things happen, naming controls in CAPS as placarded (figure 2.1-2).]
\end{nhpara}

\begin{nhsubpara}[6.8mm]{2.1.2.1}{[Key Mechanism].}
[Two lines on the single most interesting mechanism, with its rule or formula in \nowrap{one unbroken piece}.]
\end{nhsubpara}

% FIGURE 2.1-2: secondary, 44-50 mm tall. [-3.6]{50} would span y -3.6 to 46.4.
\begin{nhfigure}{46}{2.1-2}{[Title Case Caption]}{[TAB]-502A}
  \path[w] (10,10) -- (170,10);
\end{nhfigure}

\begin{nhlower}
  \begin{nhpara}{2.1.3}{PERFORMANCE AND DESIGN DATA.}\end{nhpara}
  \begin{nhtable}{2.1-1}{[Title]}
    \nhrow{[Parameter]}{[value with units]}
    \nhrow{[Parameter]}{[value with units]}
    \nhrow{[Parameter]}{[value with units]}
    \nhrow{[Parameter]}{[value with units]}
    \nhrow{[Parameter]}{[value with units]}
    \nhrow{[Parameter]}{[value with units]}
    \nhrow{[Parameter]}{[value with units]}
    \nhrow{[Parameter]}{[value with units]}
  \end{nhtable}
\nhrightcolumn
  \nhnotice{CAUTION}{[At most three lines: the one way to lose data, damage equipment or hurt someone.]}
  \begin{nhchecklist}{[PROCEDURE NAME]}{[WHO]}
    \nhstep{1}{[CONTROL] - [action]}
    \nhstep{2}{[CONTROL] - [action]}
    \nhverify{[STATE] - [expected]}
    \nhstep{3}{[CONTROL] - [action]}
  \end{nhchecklist}
\end{nhlower}

\end{document}
